\chapter{Obiectivele proiectului}\label{ch:obiective}
\pagestyle{fancy}

Sistemele de detecție a intruziunilor bazate pe învățare automată obțin, pe seturile de date de referință, rate de detecție ridicate atunci când sunt evaluate pe trafic provenit din aceeași distribuție cu cea folosită la antrenare. Această performanță nu garantează însă că sistemul rămâne robust atunci când traficul malițios este modificat în mod deliberat. Un atacator care dorește să evite detecția nu are niciun motiv să transmită atacul în forma sa „canonică”; dimpotrivă, el poate ajusta caracteristici ale traficului --- la nivel de flux sau de pachet --- astfel încât atacul să își păstreze funcționalitatea, dar să nu mai corespundă tiparelor pe care modelul le-a învățat ca fiind malițioase. Distanța dintre performanța raportată în condiții ideale și comportamentul real al sistemului în fața unor astfel de modificări constituie problema centrală abordată de această lucrare.

Tema propusă se formulează astfel: dat fiind un model de clasificare a traficului de rețea, antrenat să distingă traficul normal de mai multe tipuri de atac, se urmărește evaluarea empirică a rezistenței acestui model la variante ușor modificate ale traficului malițios. Lucrarea nu își propune construirea celui mai performant clasificator posibil, ci analiza sistematică a modului în care modificări controlate ale traficului influențează decizia unui clasificator rezonabil de bun, reprezentativ pentru soluțiile utilizate în practică. Din acest motiv, accentul cade nu pe acuratețea de detecție în sine, ci pe fragilitatea acesteia în fața perturbațiilor.

Cadrul de lucru presupune un model de amenințare (\textit{threat model}) în care atacatorul poate modifica traficul malițios generat, dar operează sub o constrângere esențială: modificările trebuie să păstreze validitatea și funcționalitatea atacului. O perturbare care ar transforma un atac într-un pachet nefuncțional nu prezintă interes practic. Prin urmare, transformările considerate se limitează la ajustări plauzibile --- de exemplu modificarea câmpului \textit{Time-to-Live} (TTL), adăugarea de octeți de umplutură (\textit{padding}) sau alterarea temporizării --- care schimbă valorile caracteristicilor extrase din trafic fără a compromite scopul atacului. Această constrângere de validitate delimitează exact domeniul temei și o diferențiază de generarea de exemple adversariale pur matematice, în care perturbarea se aplică direct în spațiul caracteristicilor, fără garanția că traficul rezultat ar putea exista într-o rețea reală.

Pe baza acestei formulări, lucrarea urmărește următoarele obiective:

\begin{description}
	\item[Obiectivul 1 --- Antrenarea unui clasificator multi-clasă de referință.] Se urmărește construirea unui model de clasificare capabil să distingă traficul normal de mai multe tipuri de atac, pe baza unui set de date public de detecție a intruziunilor. Modelul trebuie să atingă o performanță comparabilă cu cea raportată în literatură, astfel încât să constituie o bază credibilă pentru experimentele ulterioare. Evaluarea acestui model se realizează cu metrici adecvate cadrului multi-clasă dezechilibrat (precizie, \textit{recall} și scor F1 per clasă, precum și matrice de confuzie), nu doar prin acuratețe globală.

	\item[Obiectivul 2 --- Proiectarea unui modul de generare a variantelor perturbate.] Se urmărește implementarea unui mecanism care, pornind de la traficul malițios clasificat corect de model, produce variante ușor modificate ale acestuia. Modulul trebuie să permită aplicarea unor transformări de tipuri diferite (de exemplu TTL, \textit{padding}, temporizare) și pe niveluri de intensitate diferite (modificări minore, medii și accentuate), respectând constrângerea de validitate a atacului.

	\item[Obiectivul 3 --- Măsurarea ratei de evaziune.] Se urmărește determinarea proporției variantelor perturbate care reușesc să treacă nedetectate, adică să nu mai fie recunoscute drept atacuri de către model. Această măsurătoare se realizează defalcat pe tipuri de transformare și pe niveluri de intensitate, pentru a permite compararea impactului lor.

	\item[Obiectivul 4 --- Analiza sensibilității modelului.] Se urmărește identificarea tipurilor de modificări care influențează cel mai mult decizia modelului. Prin corelarea ratei de evaziune cu importanța caracteristicilor afectate de fiecare transformare, se urmărește înțelegerea motivelor pentru care anumite perturbări sunt mai eficiente decât altele. Într-un cadru multi-clasă, analiza distinge trei rezultate posibile pentru o variantă perturbată: menținerea clasificării corecte, clasificarea ca trafic normal (evaziune propriu-zisă) și clasificarea ca un alt tip de atac. Această din urmă situație oferă informații despre fragilitatea granițelor de decizie dintre clasele de atac.

	\item[Obiectivul 5 --- Dezvoltarea unei interfețe de vizualizare și testare.] Se urmărește realizarea unei interfețe simple prin care rezultatele detecției să poată fi vizualizate și prin care să se poată testa interactiv diferite niveluri de modificare a datelor. Interfața are rolul de a face rezultatele accesibile și de a permite explorarea sensibilității modelului fără a necesita intervenția directă în cod.
\end{description}

Îndeplinirea acestor obiective conduce la o evaluare de robustețe a unui sistem de detecție a intruziunilor bazat pe învățare automată, însoțită de un instrument practic de explorare. Rezultatul urmărit nu este o simplă cifră de acuratețe, ci o caracterizare a comportamentului modelului la limita dintre trafic detectat și trafic evaziv: cât de ușor poate fi indus în eroare, prin ce tipuri de modificări, și cu ce consecințe asupra clasificării. Această caracterizare oferă o perspectivă mai realistă asupra utilității unui NIDS bazat pe ML decât o oferă evaluarea clasică pe date nemodificate.